\begin{textsample}{POS Dim 3 – global\_north\_2025\_01 – Score 16.00 – global\_north\_2025\_01/120350317.txt}  \label{ex:f3_pos_262}
Harvard \textbf{University} has agreed to accept the controversial International Holocaust Remembrance Alliance ' s ( IHRA ) definition of antisemitism, after settling a lawsuit from Jewish \textbf{students} who accused other \textbf{students} of \textbf{harassment} during pro-Palestine \textbf{protests} on \textbf{campus} in 2024.
The \textbf{student} group, \textbf{Students} Against Antisemitism, filed a complaint in January 2024 in a court in Boston, saying that Harvard was failing to protect them under Title VI of the 1964 Civil Rights Act, a law that prohibits any \textbf{school} from receiving \textbf{federal} funding from discriminating against anyone based on race, colour or national origin.
The complaint cited last year ' s \textbf{pro-Palestinian} \textbf{protests} against Israel ' s war on Gaza as the source of \textbf{antisemitic} rhetoric from other \textbf{students}.
Middle East Eye has previously reported on these accusations of antisemitism against \textbf{pro-Palestinian} \textbf{protesters} at US \textbf{universities}, which often conflate support for Palestinian rights with antisemitism.
As a part of Harvard ' s settlement of the case brought by \textbf{Students} Against Antisemitism, the \textbf{university} has agreed to use the IHRA definition of antisemitism when applying its non-discrimination @ @ @ @ @ @ @ @ @ @ Dispatch Sign up to get the latest insights and analysis on Israel-Palestine, alongside Turkey Unpacked and other MEE newsletters"This resolution includes specific, meaningful actions to combat antisemitism, hate and bias on \textbf{college} \textbf{campuses} that illustrate the \textbf{University} ' s strong commitment to further protecting their Jewish and Israeli community,"Marc Kasowitz of Kasowitz Benson Torres, LLP, from the counsel for \textbf{Students} Against Antisemitism said in a statement.
The controversial IHRA definition was formulated in 2004 and published in 2005 by antisemitism expert Kenneth Stern in collaboration with other \textbf{academics} from the American Jewish Committee, a pro-Israel advocacy organisation founded at the beginning of the 20th century and based in New York.
Critics say some of the examples provided in the definition conflate antisemitism with criticism of historical policies that led to the creation of the state of Israel in 1948 - known to the Palestinians as the Nakba or"catastrophe".
The Antisemitism Awareness Act is a full frontal assault on free \textbf{speech} During the Nakba, hundreds of thousands of @ @ @ @ @ @ @ @ @ @.
Critics also say the definition conflates antisemitism with criticism of current Israeli policies against Palestinians, such as human rights abuses and the occupation by Israel.
The IHRA definition was adopted by the first Trump administration in 2019, and was reaffirmed by the Biden administration.
Harvard also said it would provide training on"combating antisemitism"and that its Office for Community Conduct would incorporate the IHRA definition when reviewing complaints of discrimination.
The move by Harvard follows other similar measures implemented by leading US \textbf{universities} to prevent any criticism of Zionism.
In August 2024, New York \textbf{University} declared Zionists a"protected \textbf{class}".
In response to Israel ' s war on Gaza that began after the Hamas-led attacks on 7 October 2023, thousands of \textbf{university} \textbf{students} across the US began organising \textbf{protests} against the war.
Many of those \textbf{protests} evolved into concerted efforts to call on their \textbf{institutions} to defund any financial stakes in companies profiting from Israel ' s war, which has killed more than 46,000 Palestinians.
In many @ @ @ @ @ @ @ @ @ @ \textbf{university} administrations, and in some cases \textbf{police} were called in to forcefully remove \textbf{protesters}.
Over the past \textbf{academic} year, \textbf{schools} have adopted more forceful measures aimed at curbing any \textbf{protests} against Israel, including issuing new \textbf{speech} guidelines and incorporating surveillance tactics against \textbf{student} organisers.
Middle East Eye delivers independent and unrivalled coverage and analysis of the Middle East, North Africa and beyond.
To learn more about republishing this content and the associated fees, please fill out this form.
More about MEE can be found here.

% matched lemmas: academic, antisemitic, campus, class, college, federal, harassment, institution, police, pro-palestinian, protest, protester, school, speech, student, university
\end{textsample}
